% Wampanoag (Wood 1634) — Source Workflow
% Issue #1367
% Status: Research complete, implementation pending
% !TEX program = xelatex
% !TEX encoding = UTF-8 Unicode

\section{Wampanoag: Wood 1634}
\label{sec:1367}

\subsection{Source Description}

\textbf{Recorder:} William Wood, English colonist living in the Massachusetts
Bay Colony (1629--1633). He returned to London to publish his observations.
Wood had no linguistic training; his orthography is purely English-based ad-hoc
spelling.

\textbf{Source:} \emph{New Englands Prospect: A true, lively, and experimentall
description of that part of America, commonly called Nevv England} (1634),
printed by Tho.\ Cotes for John Bellamie, London. Small 4to, contains a ``small
Nomenclator'' of \textasciitilde262 Massachusett words organized by semantic
category (body parts, animals, trade goods, place names, month names,
day-name classifiers).

\textbf{Language:} Wampanoag [wam].

\textbf{Corpus:} 339 records (the second-largest pre-Eliot Wampanoag source).

\textbf{Orthography:} Pre-Eliot English ad-hoc, no linguistic training. The
Eliot Massachusett writing system was still decades away. This makes Wood's
transcriptions reflect raw English-hearer perception of Wampanoag sounds,
without the regularization that Eliot later introduced.

\subsection{Recorder Biography}

William Wood was an English colonist who lived in the Massachusetts Bay Colony
from 1629 to 1633. Little is known of his life beyond his publication. He
returned to London and published \emph{New Englands Prospect} in 1634, which
went through multiple editions:

\begin{table}[htbp]
\centering
\caption{Publication History of \emph{New Englands Prospect}}
\label{tab:1367-editions}
\begin{tabular}{llll}
\hline
Edition & Year & Printer & Place \\
\hline
1st & 1634 & Tho.\ Cotes for John Bellamie & London \\
2nd & 1635 & --- & London \\
3rd & 1639 & John Dawson for John Bellamie & London \\
4th (Boston) & 1764 & Thomas and John Fleet & Boston \\
Prince Society & 1865 & --- & Boston \\
\hline
\end{tabular}
\end{table}

The 1865 Prince Society edition is the most accessible for modern researchers.
Wood's Wampanoag vocabulary appears in all editions (pp.\ 123--128 in the 1764
edition, pp.\ 94--100 in the 1865 reprint), suggesting the word list was stable
across editions with no significant revision.

\subsection{Orthography System}

Wood's orthography follows English sound-spelling conventions of the 1630s.
It is a pre-Eliot system --- the Eliot Massachusett writing system was still
decades away. Key characteristics:

\begin{itemize}
  \item \textbf{Initial /w/ retained:} \emph{Wigwam}, \emph{Web}, \emph{Wompey}
  \item \textbf{$\langle$ea$\rangle$ for /iː/:} \emph{Kean} (cf.\ Eliot's
    \emph{keen})
  \item \textbf{$\langle$s$\rangle$ for /ʃ/ inconsistently:} \emph{Sawup} where
    later sources might write \emph{sh-}
  \item \textbf{$\langle$k$\rangle$ preferred over $\langle$c$\rangle$:}
    Initial /k/ written \emph{K-} even before front vowels (\emph{Kean},
    \emph{Ksitta})
  \item \textbf{No diacritics:} Wood uses plain English letters only
  \item \textbf{Vowel length not marked:} Short vs.\ long vowel distinction
    invisible
  \item \textbf{Pidginization noted:} Goddard (1977) observes a ``good deal of
    pidginization'' in Wood's data (Salvucci 2002)
\end{itemize}

\subsection{Vowel Tables}

\begin{table}[htbp]
\centering
\caption{Wood Vowel Rules}
\label{tab:1367-vowels}
\begin{tabular}{llll}
\hline
Wood Grapheme & Likely Phoneme & Example & Notes \\
\hline
a & /a/ $\sim$ /ə/ & abbona (five), matta (no), abamocho & Short a dominant \\
e & /ɛ/ $\sim$ /i/ & aberginian (an Indian), apponees (twelve) & \\
i & /ɪ/ $\sim$ /i/ & aberginian & \\
o & /ɔ/ $\sim$ /o/ $\sim$ /ə/ & abbona, abamocho, abbomocho & Broad range \\
u & /ə/ $\sim$ /ʌ/ & sucqunnocquock (sleeps), nugge (sieve) & \\
a followed by a & /aː/ (implied) & abbona, abamacho & Double letters often for emphasis \\
oo & /uː/ $\sim$ /ɔ/ & askooke (in Winslow), week (in Wood) & \\
au & /ɔː/ $\sim$ /aw/ & pausawniscosu (half fathom), aug (I) & \\
ee & /iː/ & apponees (twelve), nees (two) & \\
ai & /eː/ $\sim$ /aj/ & pausawniscosu, appause (morning) & \\
\hline
\end{tabular}
\end{table}

\subsection{Consonant Tables}

\begin{table}[htbp]
\centering
\caption{Wood Consonant Rules}
\label{tab:1367-consonants}
\begin{tabular}{llll}
\hline
Wood Grapheme & Likely Phoneme & Example & Notes \\
\hline
b & /p/ $\sim$ /b/ & abbona (five), abamocho (devil) & Wood uses b where later sources use p \\
p & /p/ $\sim$ /b/ & apponees, appause & \\
m & /m/ & matta (no), abamocho & \\
n & /n/ & abbona, an nu ocke (a bed) & \\
t & /t/ $\sim$ /d/ & matta, abbona quit & \\
s & /s/ & apponees, pausawniscosu & \\
ch & /tʃ/ & & \\
sh & /ʃ/ & & \\
k & /k/ $\sim$ /g/ & nugge (sieve), asquoquin & \\
g & /g/ $\sim$ /k/ & nugge & \\
qu & /kʷ/ & sucqunnocquock, asquoquin & Labialized k \\
w & /w/ & naw aug (I will tell), yoaw & \\
ck & /k/ & asquoquin, an nu ocke & \\
\hline
\end{tabular}
\end{table}

\subsection{Key Orthographic Observations}

\subsubsection{b vs.\ p Alternation}

Wood writes \emph{abbona} (five) and \emph{abamocho} (devil) with \emph{b}
where the Eliot/Trumbull orthography would use \emph{p}. This shows Wood
perceived the unaspirated Wampanoag /p/ closer to English /b/ --- direct
evidence of the voicing neutralization.

\subsubsection{Geminate Consonants}

Wood frequently doubles consonants: \emph{bb}, \emph{pp}, \emph{nn}, \emph{ck},
\emph{gg}. This probably indicates the preceding vowel is short/checked rather
than true gemination.

\subsubsection{qu Digraph}

The \emph{qu} digraph consistently represents /kʷ/ --- seen in
\emph{sucqunnocquock} (sleeps), \emph{asquoquin}.

\subsubsection{Word-Initial a- Prefix}

Word-initial \emph{a-} is extremely common: \emph{abbona}, \emph{abamocho},
\emph{aberginian}, \emph{abbomocho}, \emph{abonetta}. This may be a
demonstrative/indefinite prefix that Wood transcribed as part of the root.
Roughly 25--30\% of nouns and adjectives in Wood's 339-entry word list begin
with \emph{a-}.

\subsubsection{Preaspiration /h/}

Wood does \textbf{not} write preaspirated /h/ before stops. Like Williams,
English-trained ears miss this feature. Example: \emph{matta} (no) --- modern
Wampanoag /maht-a/ or similar; \emph{sucqunnocquock} --- no \emph{h} written
despite probable /h/ in the medial cluster.

\subsection{Number System}

Wood has unusually good coverage of Wampanoag numbers (10+ distinct number
entries). This is valuable because numerals are relatively easy to verify
cross-linguistically:

\begin{table}[htbp]
\centering
\caption{Wood's Number System}
\label{tab:1367-numbers}
\begin{tabular}{lll}
\hline
Wood & Meaning & Expected Wampanoag \\
\hline
abbona & five & (p)àpàna? \\
nawhut & maybe one & nequt \\
nees & two & nees \\
quín & three & qush, shwa \\
yoâw & four & yaw, yau \\
apponees & twelve & --- \\
apponabonna & fifteen & --- \\
apponaquinta & sixteen & --- \\
apponasquoquin & nineteen & --- \\
apponenotta & seventeen & --- \\
\hline
\end{tabular}
\end{table}

\subsection{Comparison: Wood vs.\ Eliot Orthography}

\begin{table}[htbp]
\centering
\caption{Wood vs.\ Eliot Grapheme Mapping}
\label{tab:1367-vs-eliot}
\begin{tabular}{llll}
\hline
Wood & Eliot & Phoneme & Reliability \\
\hline
a & a, â & /a/, /ã/ & Low (Wood conflates length/nasal) \\
e & e & /ɛ/ & Moderate \\
i & i, u & /ɪ/, /ə/ & Low (often confused) \\
o & o, ô & /ɔ/, /ã/ & Low \\
u & u & /ə/ & Low \\
oo & oo & /uː/ & High \\
ee & ee & /iː/ & High \\
au & au, âu & /aw/, /ɔː/ & Moderate \\
b & p (never b) & /p/ & Context-dependent (normalize to p) \\
p & p & /p/ & High \\
t & t & /t/ & High \\
k/c & k & /k/ & High \\
ch & ch & /č/ & High \\
sh & sh & /š/ & High \\
qu & q & /kʷ/ & High \\
m & m & /m/ & High \\
n & n & /n/ & High \\
w & w & /w/ & High \\
y & y & /y/ & High \\
s & s & /s/ & High \\
h & h & /h/ (including preaspiration) & Missing in Wood \\
g & k (never g) & /k/ & Normalize to k \\
\hline
\end{tabular}
\end{table}

Key difference: Eliot uses a systematic vowel-length distinction (â, ô vs.\ a, o)
which Wood entirely lacks. Eliot also notes preaspiration (h) where Wood never
writes it. Eliot never uses \emph{b} or \emph{g}, always \emph{p} and \emph{k},
confirming the voicing-neutralization analysis.

\subsection{Triple-Spelling Deep Dive: ``Devil''}

The devil word is a case study in Wood's orthographic instability:

\begin{table}[htbp]
\centering
\caption{Wood's Spellings for ``Devil''}
\label{tab:1367-devil}
\begin{tabular}{ll}
\hline
Spelling & Source \\
\hline
abbomocho & Primary Wood spelling \\
abomocho & Variant (single b) \\
abamacho & Variant (a for o in second position) \\
abamocho & Variant \\
\hline
\end{tabular}
\end{table}

Proposed reconstruction: /apomaːčə/ or /apomaːčo/.

\begin{itemize}
  \item \emph{a-} prefix (demonstrative/article) $\rightarrow$ segmentable
  \item \emph{bb} $\rightarrow$ /p/ (geminate or short-vowel marker)
  \item \emph{o} $\sim$ \emph{a} in second syllable $\rightarrow$ /a/ or /ə/
    (vowel reduction)
  \item \emph{m} $\rightarrow$ /m/ (stable across all variants)
  \item \emph{ch} $\rightarrow$ /č/ (English $\langle$ch$\rangle$ = /tʃ/)
  \item final \emph{o} $\rightarrow$ /o/ or /ə/
\end{itemize}

The fact that Wood uses the same word for ``devil'' that appears in other
contexts (without Christian gloss) suggests he is mapping the Christian concept
onto an indigenous spirit-being category. This triple spelling is also useful
diagnostically: if the same word appears in Wood with multiple spellings, the
stable consonants (m, ch) are reliable while the unstable vowels (o/a, o/a)
indicate reduced/unstressed vowels that English ears struggled to categorize.

\subsection{A-Prefix Analysis}

The \emph{a-} prefix is one of the most distinctive features of Wood's
orthography. In the 339-entry word list, roughly 25--30\% of nouns and
adjectives begin with \emph{a-}. Wood treats \emph{a-} as part of the root (no
word-boundary space), suggesting either:

\begin{enumerate}
  \item \textbf{English hearer error:} Wood heard the prefixed form as a single
    phonological word and wrote it as such
  \item \textbf{Demonstrative prefix:} Wampanoag had a demonstrative/determiner
    prefix \emph{a-} (cf.\ Massachusett \emph{a-} ``that, yonder'') that was
    frequently attached to nouns in connected speech
  \item \textbf{Indefinite prefix:} Prefixed to indefinite/unspecified referents
    (a man, not the man)
  \item \textbf{Nominalizing prefix:} The \emph{a-} may be a nominalizer on verb
    stems
\end{enumerate}

Cross-source comparison: The \emph{a-} is absent in Eliot --- suggesting either
Eliot was more systematic about morphological segmentation, the \emph{a-} was
a casual-speech feature that Eliot deliberately excluded from his written
standard, or Wood heard and recorded a discourse feature (article-like prefix)
that Eliot's educated linguistic eye removed.

\subsection{Pidginization Theory}

Several scholars have suggested that Wood's word list may represent not fluent
Wampanoag but a pidginized or foreigner-talk register.

\subsubsection{Goddard's Analysis (1977, 1981)}

Ives Goddard noted that early colonial word lists often show features of
reduced morphology --- the characteristic of pidgin/foreigner-talk registers:

\begin{itemize}
  \item \textbf{Missing inflectional endings:} Wood's entries notably lack the
    inflectional suffixes that Eliot's grammar requires
  \item \textbf{Reduced pronominal marking:} Missing the complex subject/object
    marking that Eliot documents
  \item \textbf{Lexical rather than grammatical:} Wood records content words
    (nouns, verbs in citation form) rather than inflected forms
  \item \textbf{Simplified phonology:} Some consonant clusters that Eliot writes
    are simplified in Wood's transcriptions
\end{itemize}

\subsubsection{Evidence in Wood's List}

\begin{table}[htbp]
\centering
\caption{Pidginization Evidence in Wood}
\label{tab:1367-pidgin}
\begin{tabular}{lll}
\hline
Feature & Example & Pidgin Characteristic \\
\hline
Citation-form verbs & sucqunnocquock (sleeps) --- appears as-is, not inflected & Foreigner-talk citation \\
Missing possessive prefixes & an nu ocke (a bed) --- literal ``a bed'' not ``my bed'' & Reduced morphology \\
Paratactic phrases & abonetta ta sucqunnocquock (six sleeps) --- juxtaposed numerals+noun & No numeral agreement \\
Unusual word order & Some entries show English word order with Wampanoag words & Code interference \\
\hline
\end{tabular}
\end{table}

\subsubsection{Counter-Argument}

The numeral system argues against full pidginization: Wood's number series
\emph{nawhut}-\emph{nees}-\emph{quín}-\emph{yoâw}-\emph{abbona} is a close
match to the attested Wampanoag number system. Pidgins typically simplify or
restructure numeral systems; Wood's closely matches the standard language.

\subsection{Problems Encountered}

\begin{itemize}
  \item \textbf{b/p alternation:} Wood writes \emph{b} where Eliot/Trumbull use
    \emph{p} --- direct evidence of voicing neutralization. The same Wampanoag
    phoneme (unaspirated /p/) is variously written as \emph{b} or \emph{p}
    with no consistent rule.
  \item \textbf{Triple spelling:} \emph{abbomocho}/\emph{abomocho}/
    \emph{abamacho} for ``devil'' --- strongest evidence of colonial vowel
    instability in the corpus. The stable consonants (m, ch) are reliable while
    the unstable vowels (o/a, o/a) indicate reduced/unstressed vowels.
  \item \textbf{Word-initial \emph{a-} prefix:} Many entries begin with
    \emph{a-} that may be a demonstrative/indefinite prefix --- must be
    segmented before conversion. Roughly 25--30\% of nouns and adjectives are
    affected.
  \item \textbf{Geminate consonants:} Frequent doubling (\emph{bb, pp, nn, ck,
    gg}) likely marks preceding vowel as short rather than true gemination.
  \item \textbf{Preaspiration /h/ unreported:} Like all pre-Eliot sources, Wood
    does not write preaspirated /h/ before stops.
  \item \textbf{Vowel quality instability:} Multiple spellings for the same word
    show Wood had trouble distinguishing Wampanoag vowels.
  \item \textbf{Multi-word entries:} Many entries are phrases, not single
    lexemes (e.g., \emph{abonetta ta sucqunnocquock} = ``five sleeps'',
    \emph{appepes naw aug} = ``when I see it I will tell you'').
  \item \textbf{No phonetic guide:} No \emph{\textbackslash ph} entries for
    validation.
\end{itemize}

\subsection{Research Approach}

Treat Wood as unreliable but corroborative. When Wood agrees with
Trumbull/Eliot, it confirms the reading. When disagreeing, prefer Trumbull.
Normalize \emph{b} to /p/ universally. FSM with vowel averaging across Wood's
multiple spellings.

\subsubsection{Conversion Pipeline Recommendations}

\textbf{Priority 1: Normalize Consonants (High Confidence)}

\begin{itemize}
  \item \emph{b} $\rightarrow$ \emph{p} (always)
  \item \emph{g} $\rightarrow$ \emph{k} (always, but check if voiced /g/ in
    context)
  \item \emph{ck} $\rightarrow$ \emph{k} (always)
  \item \emph{qu} $\rightarrow$ \emph{kʷ} (keep as digraph or convert to q for
    phonemic)
\end{itemize}

\textbf{Priority 2: Normalize Vowels (Moderate Confidence)}

\begin{itemize}
  \item All short-vowel candidates treated as underspecified /ə/ $\sim$ /a/
    unless confirmed by Eliot
  \item Use Eliot/Trumbull/Winslow cross-reference as tiebreaker
  \item Double vowels (oo, ee) are more stable and map to /uː/, /iː/
\end{itemize}

\textbf{Priority 3: Morphological Segmentation}

\begin{itemize}
  \item Segment \emph{a-} prefix where cross-reference supports it
  \item Do not segment otherwise
  \item Flag ambiguous cases
\end{itemize}

\textbf{Priority 4: Preaspiration /h/}

\begin{itemize}
  \item Insert /h/ between vowels and following stops where Eliot confirms
  \item Default rule: when a short vowel is followed by a stop consonant in a
    stressed syllable, check Eliot for /h/
  \item Without Eliot confirmation, do NOT insert /h/ speculatively
\end{itemize}

\textbf{Priority 5: Multi-Word Entries}

\begin{itemize}
  \item Tokenize by whitespace
  \item Convert each token independently
  \item Reconstruct the phrase from converted tokens
\end{itemize}

\subsection{Conversion Workflow}
\texttt{[Pending implementation --- see Issue \#1367]}

The conversion pipeline will be implemented as a Python module in
\texttt{src/commons/parsing/}. The module will:

\begin{itemize}
  \item Treat Wood as an unreliable but corroborative source
  \item Normalize \emph{b} $\rightarrow$ /p/ universally across all Wampanoag
    entries
  \item Use FSM with vowel averaging across Wood's multiple spellings
  \item Segment \emph{a-} prefix on a word-by-word basis
  \item Handle multi-word entries by whitespace tokenization
  \item Insert preaspiration /h/ only where Eliot confirms it
  \item Verify each converted entry against expected Eliot/Trumbull forms
\end{itemize}

\subsection{Linguist Feedback Template}

\texttt{[Template --- content pending review]}

The following questions are open for linguist review:

\begin{enumerate}
  \item \textbf{b/p normalization:} Is the universal \emph{b} $\rightarrow$
    /p/ normalization correct, or are there contexts where Wood's \emph{b}
    represents a different phoneme?
  \item \textbf{a- prefix segmentation:} Is the \emph{a-} prefix a genuine
    Wampanoag morpheme (demonstrative/indefinite) or an English hearer error?
    Should it be segmented universally or on a case-by-case basis?
  \item \textbf{Pidginization:} Does the numeral system's accuracy argue against
    full pidginization, or can a language have accurate numerals alongside
    reduced morphology?
  \item \textbf{Preaspiration insertion:} Is the conservative approach (insert
    /h/ only where Eliot confirms) correct, or are there systematic contexts
    where /h/ can be reliably reconstructed?
  \item \textbf{``Devil'' reconstruction:} Is /apomaːčə/ a plausible
    reconstruction, or does the root derive from a different PA source?
\end{enumerate}

\subsection{Related Work: Colonial Recorder Perception Problem}

The core challenge of this research card --- English-speaking recorders
misperceiving and mis-transcribing Algonquian phonemes --- is a known
phenomenon documented across Eastern North America. Two researchers' work is
directly relevant:

\textbf{Dr.\ Keith Cunningham} (Linguist, Nanticoke Indian Tribe). His doctoral
dissertation \emph{A Phonological Analysis of Nanticoke With Practical
Applications for Language Revitalization} (Georgetown University) analyzes three
18th-century Nanticoke word lists (Heckewelder 1785) using the same methodology:
historical phonology from colonial recorders, with the same challenges of
multilingual informants and orthographic complexity. His 2024 Algonquian
Conference presentation \emph{A Revised Phonological Analysis of the Heckewelder
Vocabulary of Nanticoke} directly parallels the Wood analysis --- both involve
pre-missionary English recorders with no linguistic training, both show b/p
alternation as evidence of voicing neutralization, and both require PA
comparative calibration. Cunningham's work demonstrates that the colonial
recorder perception problem extends beyond SNEA into the broader Eastern
Algonquian family.

\textbf{Dr.\ Craig Kopris} (Independent scholar). His paper \emph{Les sons
wyandots perçus par des oreilles étrangères} (Wyandot sounds perceived by
foreign ears, \emph{Recherches amérindiennes au Québec}, 2014) is the exact
framing of the Wood problem: how European recorders systematically misperceived
and mis-transcribed indigenous phonemes. His \emph{Wyandot Phonology: Recovering
the Sound System of an Extinct Language} applies the same recovery methodology
to an unrelated language family (Iroquoian), demonstrating that the foreign-ear
perception problem is universal across colonial documentation of North American
languages.

The Wood corpus (339 records, the second-largest pre-Eliot Wampanoag source) is
the best case study for the pre-orthographic recorder problem. Wood's triple
spelling for a single word (\emph{abbomocho}/\emph{abomocho}/\emph{abamacho}) is
the strongest evidence in the entire SNEA corpus for colonial vowel instability
--- a phenomenon that both Cunningham (Nanticoke) and Kopris (Wyandot) document
in their respective language families.

\subsection{Key References}

\begin{itemize}
  \item Aubin, George F. 1975. \emph{A Proto-Algonquian Dictionary}. Ottawa:
    National Museums of Canada.
  \item Eliot, John. 1666. \emph{The Indian Grammar Begun}. Cambridge, MA:
    Marmaduke Johnson.
  \item Goddard, Ives. 1977. ``Some Observations on the Massachusetts
    Language.'' \emph{Papers of the 8th Algonquian Conference}.
  \item Goddard, Ives. 1981. ``Massachusetts Phonology.'' \emph{International
    Journal of American Linguistics} 47(4): 283--305.
  \item Goddard, Ives \& Bragdon, Kathleen. 1988. \emph{Native Writings in
    Massachusett}. 2 vols.\ Philadelphia: APS.
  \item Salvucci, Claudio. 2008. \emph{The Vocabulary of Massachusett}.
    Evolution Publishing.
  \item Salisbury, Neal. 1982. \emph{Manitou and Providence: Indians, Europeans,
    and the Making of New England, 1500--1643}. Oxford University Press.
  \item Trumbull, James Hammond. 1903. \emph{Natick Dictionary}. Bureau of
    American Ethnology Bulletin 25.
  \item Vaughan, Alden T. 1965. \emph{New England Frontier: Puritans and
    Indians 1620--1675}. Little, Brown.
  \item Wood, William. 1634. \emph{New Englands Prospect}. London: Tho.\ Cotes
    for John Bellamie. STC 25957.
  \item Wood, William. 1865. \emph{Wood's New England's Prospect}. Boston:
    Prince Society. (Reprint of 1634 with notes.)
\end{itemize}
